\chapter*{Working Outline}
\addcontentsline{toc}{chapter}{WorkingOutline} 

\section*{Working Titles}
\begin{itemize}
	\item Building a more effective public health practice: Rethink Public Health from a systems perspective and Transform organisational thinking, doing and learning
	\item Build a more effective public health practice: Transforming organisational thinking, doing and learning from a systems perspective
	\item Build a more effective public health practice: Transform organisational thinking, doing and learning by adopting a coherent framework of systems-based practice
	\item Build a more effective public health practice: Transform organisational thinking, doing and learning and embed a coherent framework of systems practice
	\item Build a more effective public health practice: Complexity. Systems Thinking. Critical Systems Thinking. Transform organisational thinking, doing and learning and embed a coherent framework of systems practice	
	\item Build an effective systems-based public health practice: Complexity. Systems Thinking. Critical Systems Practice. Transform organisational thinking, doing and learning.
	\item Build effective health systems thinking \& practice: Complexity. Systems Thinking. Critical Systems Practice. Transform organisational thinking, doing and learning
	\item Intractable: Why Population Health Systems Are Such a Complex Mess and What Public Health Organisations Can Do About It
	\item Project X
	\item Systems Learning Groups: Developing Organisational Capability for Systems Practice

\end{itemize}


\section*{Book Value}
\begin{itemize}
	\item \newthought{Main drivers of this book?:} 
	\begin{enumerate}
		\item I think complexity is killing effectiveness. In a world that is getting more complex, it has become difficult for the public sector to deliver and be seen to deliver improvements
		\item I have been working in the public sector since 1986 and I have never worked for an organisation that has paid more than lip service to becoming more complexity-informed in their approach to the development and delivery of public sector services
		\item I have come to the conclusion after a decade long systems thinking journey that a Systems Learning Group approach might work for public sector organisations that are led by risk-adverse leaders who cannot imagine a way forward that might break with traditional organisational management paradigms. 
	\end{enumerate}
	\item \newthought{How to start book?} I'm going to lead with a chapter on these drivers.
	\item \newthought{How to entertain/engage the reader?} Practitioner journey. I will use my journey through systems thinking to start each of the main chapters of this story and guide the reader.
	\begin{enumerate}
		\item complexity is our main challenge - my experience with pesticides
		\item we must transform thinking, doing and learning - my first public health job
		\item and adopt a cohesive systems-based practice - my most recent experiences
	\end{enumerate}
	
	\item \newthought{How to end book?} With the presentation of a framework for a cohesive systems practice that can be adopted to the needs of a team or an entire organisation's practice.
	\item \newthought{Reader value of book?}   
		\begin{enumerate}
		\item better understanding of the challenge of complexity on public health practice
		\item better understanding of how systems thinking can mitigate the challenge
		\item road map for building their own systems-based practice.
	\end{enumerate}
\end{itemize}

\section*{For each chapter:}
\begin{itemize}
	\item Hook — gets their attention, and leads to the…
	\item Thesis — your main argument, followed by your…
	\item Supporting content — which can include…
	\item Stories and examples — eventually leading to a…
	\item Recap or summary of your main points, which should reinforce the…
	\item Reader’s \#1 takeaway — which should…
	\item Segue to the next chapter — to keep the reader engaged.
\end{itemize}


\section*{Complexity Chapter}
\begin{enumerate}
	\item \newthought{Hook} My journey into complexity: In the 1980s, long before I found 'public health, I held a government internship monitoring agricultural pesticides in ground water over three summers. The first summer very few wells showed contamination, but many more wells were found to be contaminated during the second summer. By the third summer most wells were contaminated. These were drinking water wells on farms that we were sampling, and as you can imagine this was a 'problematic situation', to put it mildly. However, what you really need to know was that this wasn't the result of newly introduced pesticides or any change of practice on farms. Nor was increased rainfall to blame. Blame complexity.
	\item \newthought{Thesis} 
	\begin{itemize}
		\item (1) complexity is killing us; (2) it's largely down to our analysis (reductionist) mindset; (3) and public health's lack of diversity (do I really believe this last issue is a separate issue or just a factor in analytical thinking?)
		\item Public health is the proverbial frog in a slowly heated pot of water. It is an existential crises.
		\item Having a sound analytical framework and narrow range of thinking within the public health profession has limited our scope of action in dealing with a changing environmental situation, eroding our effectiveness...
		\item our lack of diversity: same management structures and problems, ie.'just do it' do it mentality; same education; dominated by biomedical model; where are the engineers?
	\end{itemize}

	\item \newthought{Proposition \& Evidence}
	\begin{itemize}
		\item Is our complexity challenge existential?
		\item Is analysis thinking not up to the challenge?
		\item Do we lack diversity to see problems in public health?
		\item generated list of notable public health thinkers and their challenges illustrating increasing levels of complexity
		\item surveys of what senior leaders think about complexity -- ideally of public health leaders, but probably public sector
		\item literature reviews \& my systems thinking literature database targeting complexity
		\item key studies on complexity in public health and if needed more widely in the public sector
	\end{itemize}
	
	\item \newthought{Stories, examples, case studies} Obesity Foresight project. My Dorset climate impact on older persons.
	\item \newthought{Recap}
	\item \newthought{key takeaway} Complexity is making public health less effective than it might once have been and it is seen to be less effective. Our inability to step outside our analysis paradigm is, I believe largely responsible for the situation. It has become so ingrained it is 'thinking', and we compound the ways-of-working that make it ever more difficult to escape. 
	\item \newthought{Segue to next chapter} My journey out of the analysis paradigm: Absolutely everything I looked at over the next 20 years seemed to be complex, and almost nothing in my university education equipped me to apply science and analytical skills to support the resolution of problematic situations in the real world... (new para or subsection) earth observation science... (new para or subsection) climate impacts science... (new para or subsection) academia heaven and hell... (new para or subsection) Dengue!
\end{enumerate}

\section*{Systems Thinking Chapter}
\begin{itemize}
	\item \newthought{Hook} In the late 1990s and early 2000s, long before I could articulate what systems thinking was, I had landed my first paying public health job. As a geographer with an environmental science background, I was hired to bolster the spatial and environmental capability of the NZ government's Public Health Intelligence Unit. I mapped everything, and even introduced spatial epidemiology to the public health sector, and it was very well received by colleagues. It took a while, but I eventually noticed that new insights we were getting from mapping disease and risk factors started to diminish. Geographic Information Systems, were new back then and everyone was experimenting, spatial epidemiology was a booming endeavour
	
	
	. However, the new analyses and insights weren't resolving the problematic environmental health situations that I was brought in to focus on. I was particularly stuck on one problem, and to be fair so was everyone else: Campylobacteriosis. Then I realised that I, and others, were thinking about this problem the wrong way. We need to rethink our intelligence approach.
	\item \newthought{Thesis} Adopting a new mindset, systems thinking, is more urgent than ever. There is enormous scope for increasing effectiveness, but it's not easy. There is much resistance to new ideas. Particularly when we have to reevaluate our own foundational beliefs. The resistance comes from adherents to the 'religion of science', as opposed to scientists themselves. Religion of science followers are not scientists, but people who have been convinced throughout their education and careers that science produces the only acceptable evidence on which to base decision-making and that solutions to complex evolving societal problems are best dealt with this way.
	\item \newthought{Evidence}
	\item \newthought{Stories and examples}
	\item \newthought{Recap}
	\item \newthought{Number 1 takeaway}
	\item \newthought{Segue to next chapter} In my experience, Public health was not receptive to rethinking its approach to problematic situations, they liked technical breakthroughs and they loved their numbers, but they weren't as keen on knowing why they had the problems they had. I mean they wanted to know why, of course they did, they just weren't all that keen to spend staff time and resources on thinking it through and they certainly didn't have any structured ways of doing that... (new para or subsection) a new and promising dip into academia... (new para or subsection) Exprodat...
\end{itemize}

\section*{Building Chapter}
\begin{itemize}
	\item \newthought{Hook} In the 2010s and 2020s, after trying to teach some of these ideas and seeing no opportunities to restart my public health career, I saw an opportunity to join the related field of biosecurity who seemed to welcome me with open arms. I had people around me who just 'got it'. New approaches? Just do it...  (new para or subsection) public sector consulting... (new para or subsection) a disappointing return visit to public health academia... (new para or subsection) transformation and data science in local public health intelligence... (new para or subsection) Covid and the accidental civil servant...
	\item \newthought{Thesis} A cohesively structured systems approach is long overdue and greatly needed. It hasn't happened because there is nothing cohesive about the development of the multiple systems methodologies, decision-makers by and large haven't been exposed to systems-based practice and in a world dominated by funding constraints not a lot of senior people are seriously thinking about a complete transformation of their world view. Mainstream management theory condemns us to ways-of-working developed for the industrial revolution and long past their used-by dates.  
	\item \newthought{Evidence}
	\item \newthought{Stories and examples}
	\item \newthought{Recap}
	\item \newthought{Number 1 takeaway}
	\item \newthought{Segue to next chapter} There. I've done it. I've written my book, as I've been promising myself I would do to share this journey and what I've learned with colleagues from whom have taught me so much. (new para or subsection) systems learning group journey for organisational learning...
\end{itemize}

\section*{My Questions \& Some Brief Thoughts and Reflections}
\subsection*{Systems thinking: Just the basics}
\begin{enumerate}
	\item What is 'systems thinking'?
	\begin{itemize}
		\item Why is 'systems thinking' important?
		\item Is everything with 'system' in the name actually a system?
		\item What does it mean to be a 'systems thinker'?
		\item How do I decide if I need/want to be a 'systems thinker'?
	\end{itemize}
	\item Is system dynamics the same thing as systems thinking?
	\begin{itemize}
	\item Is systems thinking a single approach?
	\item Is systems thinking the same as 'complexity informed'?	 
	\end{itemize}
	\item What else do I need to know at the beginning?
	\begin{itemize}
	\item What are systems thinking tools, I hear a lot about them?
	\item Are systems thinking models and tools the same thing?
	\end{itemize}		
\end{enumerate}

\subsection*{My Public Sector Organisation: They don't seem interested}
\begin{enumerate}
	\item Why isn't my public sector organisation using systems thinking?
	\item How do I get my organisation interested in systems thinking?
	\item How do I engage colleagues and senior leaders?
	\item I'm not having success, what do I do?
\end{enumerate}

\subsection*{Learning about systems thinking: Developing my knowledge}
\begin{enumerate}
	\item How do I develop my knowledge of systems thinking?
	\begin{itemize}
		\item How did systems thinking originate? 
		\item Who are the key thinkers in systems thinking?
		\item How much time should I devote to developing my knowledge?	
		\item How do I become a systems thinker in the public sector?
	\end{itemize}	
	\item How does a public sector workforce become systems thinkers?	
	\item How does an organisation learn?
	\begin{itemize}
		\item What does organisational learning look like?
		\item How do organisations develop a systems thinking approach?
	\end{itemize}


\end{enumerate}

\subsection*{Practising systems thinking: Developing skills and experience}
\begin{enumerate}
	\item What is a 'systems practice'?
	\item What is a 'systems practitioner'; and what do they do in a public sector organisation?
	\item How do individuals develop a systems practice?
	\item How does an organisation adopt systems practices?
	\begin{itemize}
		\item What's an organisational 'mandate'?
		\item What's a community of practice?
		\item What are organisational systems thinking structures?
	\end{itemize}
	
\end{enumerate}


\subsection*{Reflecting on the Journey: Where to now}
\begin{enumerate}
	\item Do all public sector workers need to be 'systems thinkers'?
	\item Does being a 'systems thinker' mean that we reject science and analysis thinking?
	\item What happens if we don't have more systems thinkers?	
\end{enumerate}




